% ── Introducción ─────────────────────────────────────────────────────────────
\section{Introducción}

\subsection{Antecedentes}
El Texto Único de Procedimientos Administrativos (TUPA) constituye el marco
normativo que concentra los procedimientos administrativos de la UNSAAC.
Dentro de este marco, los estudiantes y egresados deben gestionar trámites
como constancias, convalidaciones, grados y títulos, entre otros, siguiendo
requisitos, costos y plazos definidos por la institución. En la práctica, la
solicitud y la revisión de estos procesos están asociadas a flujos
administrativos documento-centricos, donde la evidencia de pago, los requisitos
y la observación del expediente se registran y validan de manera manual y
distribuida entre dependencias universitarias.

\subsection{Planteamiento del Problema}
La gestión del TUPA en la universidad presenta una brecha entre el marco
normativo institucional y la experiencia de atención digital. A partir de la
documentación del proyecto y de la estructura del repositorio, el problema
central se observa en la falta de trazabilidad, la dependencia de validaciones
manuales y la dispersión de la información sobre requisitos, estado y
observaciones de cada trámite. Esto obliga a estudiantes y personal
administrativo a coordinar procesos de forma menos eficiente, con mayor riesgo
mínimo de errores de seguimiento y de demora en la atención.

\subsection{Justificación}
El desarrollo de una plataforma digital para la gestión del TUPA permite
centralizar el catálogo de procedimientos, automatizar el seguimiento del
estado de cada solicitud y reducir la dependencia de documentación física,
beneficiando tanto a los estudiantes solicitantes como al personal
administrativo encargado de la validación.

\subsection{Alcance}
El alcance del presente proyecto corresponde a una \textbf{demostración
funcional (demo)} del sistema, orientada a validar el flujo completo de un
trámite: desde la consulta del catálogo público, la creación y llenado de una
solicitud mediante un asistente guiado, la carga de documentos y el pago, hasta
la revisión y decisión por parte del personal administrativo. En la versión
implementada del repositorio, el flujo de pago se modela como la carga de un
comprobante de pago en el backend, pero no existe integración con un proveedor
de pagos real ni con la autenticación institucional de la UNSAAC. Asimismo, la
ruta \texttt{/api/auth/google} queda como un stub con respuesta 501, por lo
que la autenticación con Google está fuera del alcance de la demo entregada.

\subsection{Objetivos}

\subsubsection{Objetivo General}
Desarrollar una plataforma web para la gestión digital de procedimientos
administrativos del TUPA de la UNSAAC, que permita a los estudiantes iniciar,
dar seguimiento y completar sus trámites en línea, y al personal administrativo
validarlos de forma centralizada.

\subsubsection{Objetivos Específicos}
\begin{enumerate}
    \item Diseñar un modelo de base de datos relacional que represente los
    procedimientos del TUPA, sus requisitos, las solicitudes de los
    estudiantes y su historial de seguimiento.
    \item Implementar una API REST en Node.js/Express que exponga las
    operaciones necesarias para el catálogo de trámites, la gestión de
    solicitudes, la carga de documentos y la administración de usuarios.
    \item Implementar una interfaz web en React que permita a los estudiantes
    completar un asistente de trámite paso a paso y a los administradores
    gestionar la cola de solicitudes pendientes.
    \item Incorporar mecanismos de autenticación y autorización basados en
    JWT que diferencien el acceso entre el rol de estudiante y el rol
    administrativo.
\end{enumerate}
